\section{Selves}

This is the most speculative material in the paper, and it begins with a firm caveat. Raw vibes are not physics. The ternary tones are values, not particles, and the substrate does not map onto the world in one step. Between the base and any physics there are one to three emergent middle layers, and selves, where they appear at all, live in that middle and not in the base. Everything below is labeled by tier, and where a claim is conjecture it is marked \textbf{open}, Tier O.

The upper program rests on one constraint that doubles as a prediction. A closed, deterministic, reversible base sits at the causal-reduction end, where the finest micro account stays the best account and no coarse summary improves on it. Causal emergence, where a coarse layer becomes a sharper cause than the cells beneath it, needs noise or degeneracy to average away, and a reversible rule has neither. So autonomy can enter only through one of two doorways: lossy coarse-graining, where a many-to-one projection injects effective noise, or open subsystems trading tone with the ever-growing wake. The yardstick for whether a layer is causally real is effective information, how much intervening on a cause pins down the effect \cite{hoel2013}. This much is Tier F.

\begin{finding}
Selves cannot come from the closed reversible base. They can arise only in the lossy middle. The pure base churns and makes no bound body, which is the recorded negative the constraint predicts.
\end{finding}

The closed-base case was run and reported. At over a million docks, with charge and momentum conserved exactly and the run bit-identical on two machines, a proto-self does not bind. It spreads ballistically at about 1.3 cells per beat. The bare rule delivers identity and radiation, not a bound body. This is the closed base behaving exactly as the constraint says it must, a measured negative tracked in the selves arena, the E-SLF experiments under the decisive labels E-DEC-02 and E-DEC-06. What stays open is that a bound, self-repairing body needs one more discrete ingredient, an attraction, which the bare atoms do not supply. Tier O.

The right test for a candidate self is a body-hit: perturb the structure's own charges and watch its body. Binding has to be confirmed by spatial localization, a pattern that stays put and relaxes back when pushed, and not by a recurring scalar, because a recurring number is exactly what a simple gas already produces without binding anything. A real self should register as an irreducible high-order cause whose interventions carry more effective information than its own cells do. The decisive run computes that quantity in three cases, the closed base, a lossy coarse-graining, and an open subsystem, expecting selves to appear only in the last two. This is the E-SLF arena, with the controls that the closed base shows reduction and a scramble destroys the structure.

The middle is also where freedom could appear. The claim here is deliberately limited. Two things both called open must be kept apart. Reality-openness asks whether the same total state could yield two different futures, and on a deterministic reversible rule the answer is no, always. Agent-openness asks whether a self holds a picture of alternatives and cannot read its own next act until it settles, and there the answer can be yes. Freedom, on this reading, is the product of modeled openness and self-unreadability, a quantity that is zero for a bare tone with nothing to model, zero for the all-containing whole with nothing outside to model, and large for a part modeling a world bigger than its view. Determinism holds underneath, unchanged. This framing draws on Friston's Markov blanket read as a Lyapunov-settling agent with no literal inference \cite{friston2010} and on the self as a self-maintaining predicting boundary \cite{fields2018}. The framing of laws as which transformations are possible and impossible is the admissibility lens \cite{deutsch2015, marletto2017}. All of this is Tier O, the interpretive layer, stated and clearly labeled rather than derived.

The measurable parts of this section are firm. The deepest claim remains a conjecture, Tier O. Conservation is the precondition that lets a pattern settle into an attractor and persist, which is necessary for a self but not yet sufficient, since it does not by itself pull the pattern into a bound body. The model does not prove free will. It defines freedom as two measurable quantities, shows the product should peak in the middle for reasons that trace to the firm base, and predicts the one place a living self could ever come from while ruling out the place it cannot. The identification of the tones with felt pain, peace, and pleasure is the monist reading. It is the one place the model reaches past physics, and it is offered tentatively, as the base axiom seen from the inside rather than a derived result.
